\section{O que anda junto com o quê}
\label{sec:correlacao}

A Figura~\ref{fig:correlacao} resume, num único mapa, o quanto cada par
de indicadores anda junto. A medida usada (correlação de Spearman) vai
de 0, quando dois indicadores não têm relação, a 1, quando crescem
sempre juntos. A mortalidade ajustada foi retirada do mapa por ser
estatisticamente idêntica à mortalidade geral (correlação de 1,000
após arredondamento entre as duas; todas as demais relações coincidem
casa a casa).

Os pares mais fortes fazem sentido imediato para a gestão:
\textbf{faturamento por saída e fração de alta complexidade (0,77)}, faturamento e
tempo de permanência (0,52), e as duas ocupações entre si (0,54). A
\textbf{complexidade estrutural conversa com tudo}: 0,56 com alta
complexidade, 0,62 com ocupação de internação, 0,58 com faturamento, 0,60 com
TMP e 0,25 com mortalidade, o retrato numérico do papel de grande
divisor que a Seção~\ref{sec:complexidadeporte} descreveu. Já a
mortalidade se associa mais à ocupação de UTI (0,36) do que à de
internação (0,22), na direção esperada: onde a casuística é mais grave,
usa se mais UTI e morre se mais.

\comoler{procure a linha de um indicador e percorra as colunas: cada
célula diz o quanto aquele par anda junto (0 sem relação, 1 relação
perfeita). A diagonal vale sempre 1 porque compara o indicador com ele
mesmo. Repare na coluna da complexidade estrutural, escura quase
inteira: é o confundidor que precisa ser descontado em qualquer
comparação.}

\begin{figure}[H]
\centering
\includegraphics[width=.62\textwidth]{fig_ae_10_correlacao_spearman.png}
\caption{Mapa de correlações entre os indicadores, as ocupações e a
complexidade estrutural. Azul mais escuro indica relação mais forte;
todos os valores estão impressos.}
\label{fig:correlacao}
\end{figure}

\comofeito{mapa de calor. Nos dois eixos aparecem os mesmos itens --- os
sete indicadores, as duas taxas de ocupação e a complexidade estrutural
--- e cada célula traz a \textbf{correlação de Spearman} entre o par
correspondente, medida que vai de 0 (sem relação de ordem) a 1 (crescem
sempre juntos). Azul mais escuro indica relação mais forte; a diagonal
vale sempre 1 porque compara cada item consigo mesmo. Escolheu-se
Spearman por captar relações de ordem sem exigir forma linear.}

A Figura~\ref{fig:variancia} e a Tabela~\ref{tab:variancia} respondem a
outra pergunta essencial: \textbf{quando um indicador varia, essa
variação vem de comparar hospitais diferentes ou do mesmo hospital
mudando ao longo dos anos?} A resposta é contundente: a variação vem
quase toda das diferenças entre hospitais. Na fração de alta complexidade,
\textbf{apenas 4,5\% da variação total acontece dentro do mesmo hospital}
ao longo da década; no faturamento são 8,3\%, no tempo de permanência
13,1\%, na mortalidade 16,0\%, na ocupação de internação 16,7\%, e só a
ocupação de UTI passa de um terço (38,2\%).
Em termos de gestão: \textbf{cada hospital é muito parecido consigo
mesmo, ano após ano, e muito diferente dos vizinhos}. Os indicadores
têm forte inércia institucional, e mudanças reais dentro de um hospital
são eventos raros e valiosos de observar, o que dá ainda mais peso aos
cinco casos de conversão da Seção~\ref{sec:conversores}.

\begin{table}[H]
\centering
\caption{De onde vem a variação de cada indicador: parcela que ocorre
dentro do mesmo hospital ao longo dos anos e parcela que vem das
diferenças entre hospitais.}
\label{tab:variancia}
\small
\begin{tabular}{lrr}
\toprule
Indicador & Dentro do hospital (\%) & Entre hospitais (\%) \\
\midrule
Fração alta complex.   & 4,5  & 95,5 \\
Faturamento por saída        & 8,3  & 91,7 \\
TMP                    & 13,1 & 86,9 \\
Mortalidade geral      & 16,0 & 84,0 \\
Mortalidade ajustada   & 16,0 & 84,0 \\
Ocupação internação    & 16,7 & 83,3 \\
Ocupação UTI           & 38,2 & 61,8 \\
\bottomrule
\end{tabular}
\end{table}

\comofeito{cada linha é um indicador e as duas colunas repartem 100\% da
sua variação total em duas parcelas: a que ocorre \textbf{dentro do mesmo
hospital} ao longo dos anos e a que vem das \textbf{diferenças entre
hospitais}. As parcelas saem da decomposição da variância de cada
indicador nessas duas componentes.}

\comoler{quanto menor a barra, mais estável o indicador dentro de cada
hospital e mais as diferenças observadas são traço permanente da
instituição. A fração de alta complexidade, com 5\%, é quase uma
impressão digital do hospital; a ocupação de UTI, com 38\%, é a que mais
responde ao tempo, em boa parte por causa da pandemia.}

\begin{figure}[H]
\centering
\includegraphics[width=.62\textwidth]{fig_ae_14_variancia_within.png}
\caption{Parcela da variação de cada indicador que acontece dentro do
mesmo hospital ao longo do tempo.}
\label{fig:variancia}
\end{figure}

\comofeito{gráfico de barras horizontais. Cada barra é um indicador
(eixo vertical) e o comprimento mostra a \textbf{parcela da variação
total daquele indicador que ocorre dentro do mesmo hospital ao longo dos
anos} (eixo horizontal, em \%); o restante, até 100\%, vem das
diferenças entre hospitais. A parcela foi obtida decompondo a variância
de cada indicador em uma componente dentro do hospital e outra entre
hospitais.}

